% !TEX program = xelatex
% Bend 2.0.36 / 2026-10-16; reproducible source of the accompanying PDF.
\documentclass[11pt,a4paper,openany]{book}
\usepackage[margin=22mm,top=25mm,bottom=24mm,headheight=16pt]{geometry}
\usepackage{fontspec}
\usepackage{xeCJK}
\setmainfont{Noto Serif}
\setsansfont{Noto Sans}
\setmonofont{Noto Sans Mono CJK JP}[Scale=0.84]
\setCJKmainfont{Noto Serif CJK JP}
\setCJKsansfont{Noto Sans CJK JP}
\setCJKmonofont{Noto Sans Mono CJK JP}
\usepackage{amsmath,amssymb,mathtools,booktabs,tabularx,longtable,multicol}
\usepackage{ragged2e}
\newcolumntype{L}[1]{>{\RaggedRight\arraybackslash}p{#1}}
\usepackage{tikz}
\usetikzlibrary{arrows.meta,positioning,calc,fit,shapes.geometric}
\usepackage[most]{tcolorbox}
\usepackage{listings}
\usepackage{xcolor}
\usepackage{graphicx}
\usepackage{enumitem}
\usepackage{fancyhdr}
\usepackage{hyperref}
\usepackage{microtype}
\usepackage{needspace}
\usepackage{etoc}
\usepackage{float}
\definecolor{ink}{HTML}{17263C}
\definecolor{blue}{HTML}{176A9A}
\definecolor{teal}{HTML}{167A72}
\definecolor{warm}{HTML}{C16D25}
\definecolor{paper}{HTML}{F5F7FA}
\definecolor{soft}{HTML}{E6EDF2}
\definecolor{codebg}{HTML}{F4F7FA}
\definecolor{muted}{HTML}{536378}
\hypersetup{colorlinks=true,linkcolor=blue,urlcolor=teal,pdftitle={ライフゲームで学ぶ はじめてのBend2},pdfauthor={Nasu726},pdfsubject={Bend 2 hands-on programming and Game of Life}}
\setlength{\parskip}{.5em}
\setlength{\parindent}{0pt}
\raggedbottom
\setlist{nosep,leftmargin=2em}
\fancyhf{}
\fancyhead[LE,RO]{\sffamily\small\textcolor{muted}{ライフゲームで学ぶ はじめてのBend2}}
\fancyhead[LO,RE]{\sffamily\small\textcolor{muted}{\nouppercase{\leftmark}}}
\fancyfoot[C]{\sffamily\small\thepage}
\pagestyle{fancy}
\renewcommand{\headrulewidth}{0.3pt}
\usepackage{titlesec}
% 標準bookクラスのpartは独立した見出しページを作ります。
% straightクラスへ変え、見出しの直後に当該Partの目次を配置します。
\titleclass{\part}{straight}
\titleformat{\part}[display]
  {\sffamily\bfseries\color{ink}\raggedright}
  {\color{teal}\fontsize{16}{21}\selectfont PART \thepart}
  {5mm}
  {\fontsize{27}{35}\selectfont}
\titlespacing*{\part}{0pt}{30mm}{15mm}
\titleformat{\chapter}[display]{\normalfont\sffamily\bfseries\color{ink}}{\fontsize{12}{14}\selectfont Chapter \thechapter}{10pt}{\fontsize{24}{30}\selectfont}
\titleformat{\section}{\Large\sffamily\bfseries\color{ink}}{\thesection}{1em}{}
\titleformat{\subsection}{\large\sffamily\bfseries\color{blue}}{\thesubsection}{1em}{}
\lstdefinelanguage{BendTwo}{
  morekeywords={import,def,law,type,is,match,case,do,return,for,exs,where,unsafe},
  sensitive=true,
  morecomment=[l]{\#},
  morestring=[b]",
}
\lstset{language=BendTwo,basicstyle=\ttfamily\footnotesize,
 keywordstyle=\color{blue}\bfseries,commentstyle=\color{muted},
 stringstyle=\color{teal},numbers=left,numberstyle=\tiny\color{muted},
 stepnumber=1,numbersep=8pt,breaklines=true,breakatwhitespace=true,
 showstringspaces=false,columns=fullflexible,keepspaces=true,upquote=true,
 tabsize=2,frame=single,rulecolor=\color{soft},backgroundcolor=\color{codebg},
 xleftmargin=12pt,xrightmargin=3pt,framexleftmargin=12pt,
 aboveskip=9pt,belowskip=9pt,captionpos=b}
\newtcolorbox{pointbox}[1][]{enhanced,breakable,colback=paper,colframe=blue!32,
 boxrule=.6pt,arc=2mm,left=4mm,right=4mm,top=3mm,bottom=3mm,
 title={\sffamily\bfseries 要点},coltitle=ink,fonttitle=\sffamily\bfseries,#1}
\newtcolorbox{exercisebox}[1][]{enhanced,breakable,colback=white,colframe=teal!55,
 boxrule=.7pt,arc=1.5mm,left=4mm,right=4mm,top=3mm,bottom=3mm,
 title={\sffamily\bfseries ハンズオン／確認問題},coltitle=ink,#1}
\newtcolorbox{cautionbox}[1][]{enhanced,breakable,colback=warm!5,colframe=warm!65,
 boxrule=.7pt,arc=1.5mm,left=4mm,right=4mm,top=3mm,bottom=3mm,
 title={\sffamily\bfseries 注意},coltitle=ink,#1}
\newcommand{\code}[1]{\texttt{#1}}
\newcommand{\file}[1]{\textsf{\color{blue}#1}}
\newcommand{\term}[1]{\textbf{#1}}
\newcommand{\key}[1]{\textcolor{blue}{\textbf{#1}}}
\newcommand{\readcode}[2]{\lstinputlisting[caption={#2}]{src/#1}}
\newcommand{\exercise}[1]{\begin{exercisebox}#1\end{exercisebox}}
% Part別の章一覧は .toc から自動抽出。章を増減しても同期します。
\newcommand{\partchaptertoc}{%
  \begingroup
  \etocsetnexttocdepth{chapter}%
  \etocsettocstyle{%
    \vspace{2mm}\begin{tcolorbox}[colback=paper,colframe=blue!32,boxrule=.5pt,arc=2mm,
        title={\sffamily\bfseries このPartの目次},coltitle=ink,fonttitle=\sffamily\bfseries]
    \small\sffamily
  }{\end{tcolorbox}}%
  \localtableofcontents
  \endgroup
}

\begin{document}
\frontmatter
\begin{titlepage}
  \pagecolor{paper}
  \begin{tikzpicture}[remember picture,overlay]
    \fill[blue] ([xshift=-2.2cm,yshift=-8.3cm]current page.north east) circle (4.0cm);
    \fill[teal!70] ([xshift=-2.2cm,yshift=-8.3cm]current page.north east) circle (2.6cm);
    \foreach \x/\y in {0/0,1/0,2/0,2/1,1/2,0/2,3/3,4/2,5/1,4/0} {
      \fill[white,opacity=.55] ([xshift=-4.3cm,yshift=-10.7cm]current page.north east) ++(\x*.37cm,\y*.37cm) rectangle ++(.28cm,.28cm);
    }
  \end{tikzpicture}
  \vspace*{34mm}
  {\sffamily\fontsize{15}{19}\selectfont\color{blue} 小さなプログラムから一つの作品へ}\par
  \vspace{13mm}
  {\sffamily\bfseries\fontsize{22}{29}\selectfont\color{blue} ライフゲームで学ぶ\par}
  \vspace{5mm}
  {\sffamily\bfseries\fontsize{37}{48}\selectfont\color{ink} はじめてのBend2\par}
  \vspace{12mm}
  {\sffamily\Large Bend 2の文法・並列計算・証明を、ライフゲームで学ぶ}\par
  \vfill
  {\sffamily\large 本文・演習・Bendソース・独立テスト付き}\par
  \vspace{5mm}
  {\sffamily\large 第1版 \quad / \quad Bend 2 基本文法}\par
  \vspace{15mm}
\end{titlepage}
\nopagecolor
\chapter*{本書の読み方と検証状態}
\addcontentsline{toc}{chapter}{本書の読み方と検証状態}
本書で作るものは、\term{8×8のトーラス型ライフゲーム（B3/S23）を5世代表示するターミナル・シミュレータ}です。
完成版は\file{src/08\_life\_terminal.bend}に収録しています。ただし、初めから完成版を丸写しするのではありません。
まずBend 2の型・関数・配列などを学びます。その後ライフゲームの規則を確認し、1セルの更新、盤面と近傍の計算、世代更新、文字表示を順に組み合わせます。

最初にBend 2の値・型・関数などを学び、次にライフゲームを作り、CPU/GPU並列実行と形式的証明を学びます。各章では「出力を予想する $\rightarrow$ 実行する $\rightarrow$ 改造する $\rightarrow$ 結果を確認する」という順序で進めてみましょう。

\begin{pointbox}
\textbf{付属テストの結果}：Bend 2.0.37によるサンプルの型検査に成功しています。ライフゲームの生死規則18ケースと5世代の盤面も、別途作成したPythonの参照実装と照合済みです。\textbf{テストは、調べた入力について正しい結果が得られたことを示すもの}です。
\end{pointbox}

各章の演習では、まず自分で考えてから付属のサンプルを実行してみましょう。分からない用語は直前の解説や用語集を参照できます。

\section*{同梱ファイル}
\begin{tabularx}{\linewidth}{@{}L{65mm}X@{}}
\toprule
ファイル & 内容\\
\midrule
\file{src/*.bend} & 型・関数・配列などの入門用例と、ライフゲーム本体、並列呼び出しと証明の例。\\
\path{src/12_affine_probe_INVALID.bend} & 意図的に型エラーを起こす練習用コード。\\
\file{scripts/reference.py} & 期待する盤面を独立に生成するPythonプログラム。\\
\path{scripts/test_bend_runtime.py} & Bendコードを実行し、出力の不一致や実行失敗を検出するテスト。\\
\file{tests/} & ライフゲームの参照テストと教材の構成検査。\\
\file{Bend2\_Life\_Textbook.tex} & 本文のLaTeX原稿。\\
\bottomrule
\end{tabularx}


\chapter*{学習前に知っておきたい基本用語}
\addcontentsline{toc}{chapter}{学習前に知っておきたい基本用語}
プログラミングを始めたばかりの方も、本書を読み進められるようにします。
ここでは、まず\term{プログラム}が「コンピュータに行わせる処理を書いたもの」であることを確認します。
そのプログラムを人が読んだり編集したりできる文字列が\term{ソースコード（source code）}です。
難しい言葉をすべて覚える必要はありません。コードに出てきたときに、このページへ戻ってください。

\begin{longtable}{@{}L{47mm}p{103mm}@{}}
\toprule 用語 & 意味と本書での使い方\\\midrule\endhead
\term{値（value）} & 数値の\code{3}や文字列の\code{"Life"}など、計算の対象になるものです。\\
\term{変数（variable）} & 値に付ける名前です。\code{x}や\code{board}などを使います。Bend 2では利用回数に制約がある場合があります。\\
\term{式（expression）} & 計算して値になる記述です。たとえば\code{(2 + 3 : U32)}です。\\
\term{関数（function）} & 入力を受け取り結果を返す処理です。\code{calculate}は整数の計算を行います。\\
\term{引数（argument）} & 関数に渡す具体的な値です。\code{calculate(2,3)}の\code{2}と\code{3}が引数です。関数を宣言するときの\code{x}などの名前は\term{仮引数（parameter）}と呼びます。\\
\term{返り値（return value）} & 関数が計算の結果として返す値です。\code{calculate(2,3)}の返り値は\code{13}です。\\
\term{型（type）} & 値の種類と、可能な操作を表す情報です。\code{U32}は32ビット符号なし整数、\code{String}は文字列の型です。\\
\term{配列（array）} & 複数の値を順番に並べ、位置を指定してアクセスする構造です。64要素に固定された仕組みではありません。\\
\term{添字（index）} & 配列の要素がどの位置にあるかを示す番号です。最初の位置を0として数えます。\\
\term{再帰（recursion）} & 関数が自分自身を呼ぶ処理です。終わる条件を設ける必要があります。\\
\term{副作用（side effect）} & 値を返すだけでなく、画面表示やファイル書き込みなど外部へ働きかけることです。\\
\term{コンパイル（compilation）} & ソースコードを、実行できる形式などへ変換する工程です。\term{型検査（type checking）}は、各式の型が合っているかを調べる検査です。\\
\term{並列計算（parallel computation）} & 独立した計算を複数同時に進めることです。\\
\term{形式的証明（formal proof）} & 法則を対象のすべての値について示し、その根拠を処理系で検査できる形にしたものです。\\
\term{テスト（test）} & 期待した入力と出力を用意し、実際の動作が一致するかを調べることです。すべての入力に対する数学的な証明とは異なります。\\
\bottomrule
\end{longtable}

\begin{pointbox}
Bendの型名や構文（\code{U32}、\code{Data}、\code{IO}など）は、英語の綴りのままコードに書きます。
\end{pointbox}

\tableofcontents
\mainmatter
\clearpage
\part{Bend 2の基本文法を学ぶ}
\partchaptertoc
\begin{pointbox}
Bend 2の値と型を理解したら、関数、場合分け、配列の操作、リスト、高階関数、クロージャ、モナド、モジュールへ順に進みます。小さなプログラムを動かしながら、構文と結果を確認してみましょう。
\end{pointbox}
\chapter{環境構築と最初の実行}
\begin{tcolorbox}[colback=blue!3,colframe=blue!16,boxrule=.35pt,arc=1mm,left=3mm,right=3mm,top=2mm,bottom=2mm]
\small 本章では、\term{ターミナル（terminal）}という文字入力の画面からBendを起動します。コマンドを一つ実行し、表示が変わるところまで進みましょう。
\end{tcolorbox}

\section{導入前に確認すること}
本書ではBend 2を使用します。旧Bend 1と構文が異なるため、インストールしたコンパイラの\code{bend version}を確認しましょう。WindowsではWSLのLinux環境を利用できます。macOSやLinuxではターミナルから操作します。

\begin{lstlisting}[language={},caption={導入、動作確認、一次資料の表示}]
curl -fsSL https://bend-lang.com/install.sh -o install-bend.sh
# Inspect install-bend.sh before execution.
sh install-bend.sh
export PATH="$HOME/.bend/bin:$PATH"
bend version
bend guide
bend base
\end{lstlisting}
\code{bend version}の値と表示されたパスを\file{VERSION\_LOG.md}へ記録します。インストール用シェルを内容確認せず実行しません。\code{bend update}がある環境では勝手に更新した後の資料と旧コードを混ぜず、バージョンを記録してから学習します。

\section{Bendのプログラムを実行する}
まず、画面に短い文章を表示するプログラムを動かします。次の例を見てください。分からない記号があっても、今はすべて暗記する必要はありません。
\readcode{01_hello.bend}{はじめてのBend 2（01\_hello.bend）}

このコードが行うのは、\textbf{画面に「Hello, Bend 2」と表示すること}だけです。上から一行ずつ見ていきましょう。
\begin{enumerate}
\item \code{import Base}は、Bendに用意されている基本的な機能を使えるようにします。ここでは文字を表示するために使います。
\item \code{def main() -> IO(Unit):}は、プログラムの開始点となる\code{main}という名前の処理を定義します。\code{def}は処理を定義するための言葉です。\code{->}以降は処理の結果に関する型の指定ですが、詳しい読み方は後の章で学びます。
\item \code{do IO<Unit>:}は、画面への表示などを順に実行する部分の始まりです。前の行より右に字下げされた行が、この部分に含まれます。
\item \code{IO.print("Hello, Bend 2")}が、括弧の中の文章を画面に表示します。ダブルクォーテーション\code{"}で囲んだ部分が表示する文字列です。
\end{enumerate}

\term{IO（入出力）}とは、画面への表示やキーボードからの入力など、プログラムの外とのやり取りです。\code{Unit}は「特に返す値がない」ことを表す型です。どちらも後の章で詳しく学びます。

ターミナルで、このリポジトリのフォルダーを開き、次の一行を入力してください。
\begin{lstlisting}[language={},caption={最初のプログラムを実行する}]
bend src/01_hello.bend
\end{lstlisting}
\code{bend}はプログラムを動かすコマンド、\code{src/01\_hello.bend}は実行するファイルの場所です。成功すると、次のように表示されます。
\begin{lstlisting}[language={},caption={画面に表示される内容}]
Hello, Bend 2
\end{lstlisting}

ここまでできたら、\file{src/01\_hello.bend}の\code{"Hello, Bend 2"}を好きな文章に変えて保存し、もう一度同じコマンドを実行してみましょう。表示が変われば成功です。

\begin{pointbox}
実行せずにプログラムの書き方を確認したいときは、\code{bend src/01\_hello.bend --check-only}を使えます。ネイティブ実行ファイルを作る方法もありますが、最初は上の\code{bend}コマンドだけで十分です。
\end{pointbox}

\section{型エラーの表示を確認する}
\code{IO.print}の文字列を数値に変えてみましょう。次に型注釈を消してみます。さらに返り値の型を\code{U32}にして表示命令を残し、どのエラーが出るか観察します。エラーメッセージのうち\term{expected}と\term{observed}がどの式を指しているかを説明します。

\exercise{(1) \code{Hello, Bend 2}の文字列を変更して実行してみましょう。(2) \code{bend guide}で「Types」に関する項目を探してみましょう。(3) 導入版・OS・CPU・Clangの情報を\file{VERSION\_LOG.md}に保存してみましょう。\textbf{合格条件}: ソースを編集して標準出力を変えられること。}

\chapter{値とデータ型を理解する}
\begin{tcolorbox}[colback=blue!3,colframe=blue!16,boxrule=.35pt,arc=1mm,left=3mm,right=3mm,top=2mm,bottom=2mm]
\small プログラムを書く前に、\term{値（value）}と\term{型（type）}を区別しましょう。Bend 2の型は「どんな値を表し、何をしてよいか」を決めます。
\end{tcolorbox}

\section{値と型の対応を読む}
\term{データ型（data type）}とは、値の種類を定めるものです。例えば\code{42}は数値、\code{"hello"}は文字列です。Bend 2で頻出する型を先に見ておきます。

\begin{tabularx}{\linewidth}{@{}L{23mm}L{42mm}X@{}}
\toprule 型 & 値の例 & 用途\\\midrule
\code{Nat} & \code{3n} & 0以上の自然数を表します。\\
\code{U32} & \code{42} & 0から$2^{32}-1$までの符号なし整数です。\\
\code{Bool} & \code{True\{\}} & 真または偽を表します。\\
\code{String} & \code{"hello"} & 文字列を表します。\\
\code{F32} & 32ビット浮動小数点 & 小数を近似して扱う数値型です。\\
\bottomrule
\end{tabularx}
同じ数字に見えても\code{3n}と\code{3}は異なる型の値です。Bend 2は型の暗黙変換が少ないため、型を意識してコードを読む必要があります。\term{型注釈（type annotation）}は、値や式の型を明示する記法です。例えば\code{(3 + 4 : U32)}は、\code{U32}型の足し算を表します。

\section{選択肢を新しい型として定義する}
既存の型だけでは表しにくい値は、\term{ユーザー定義型（user-defined type）}を作って表現できます。次の例では、\term{コンストラクタ（constructor）}の\code{Idle\{\}}と\code{Busy\{\}}という二種類だけの状態を定義します。コンストラクタとは、\textbf{その型の値を組み立てるための名前}です。
\readcode{28_data_types.bend}{二つの状態を持つデータ型（28\_data\_types.bend）}
最初の\code{type Status is Data:}が型の宣言、続く二行が作れる値の候補です。最後の\code{main}はプログラムの入口で、詳しい関数の書き方は次章で学びます。ここでは、\code{Idle\{\}}が\code{Status}型の値になることを読み取りましょう。
\begin{pointbox}
\term{代数的データ型（algebraic data type; ADT）}は、候補を列挙したり、候補ごとに別のデータを持たせたりできる型です。例えば\code{Some\{value\}}と\code{None\{\}}で「値がある／ない」を表す\code{Maybe}もその一例です。\code{is Data}は値を再利用できる分類を表しますが、所有権とあわせて後の章で扱います。
\end{pointbox}
\exercise{\code{Status}に\code{Paused\{\}}を追加してみましょう。\code{Idle\{\}}と\code{Busy\{\}}以外の状態を、型を変えずに作れるか考えてみましょう。}

\chapter{関数を定義して数値を計算する}
\begin{tcolorbox}[colback=blue!3,colframe=blue!16,boxrule=.35pt,arc=1mm,left=3mm,right=3mm,top=2mm,bottom=2mm]
\small \term{関数（function）}、\term{引数（argument）}、\term{返り値（return value）}を短い数値計算で確認します。特定のアプリケーションに依存しない例で、型注釈と式の読み方を身につけましょう。
\end{tcolorbox}

\section{引数と返り値の型}
Bendでは\code{def f(x: T) -> R:}の形で関数を宣言します。戻り値は最後の式であり、Pythonのように常に\code{return}が必要なのではありません。\code{U32}は0以上の整数を32ビットに格納する型です。\code{Nat}は0、1、2、……という\term{自然数（natural number）}の型で、再帰や型の添字に使います。\code{3n}は\code{Nat}のリテラルで、\code{3}とは扱いを区別します。

\readcode{02_arithmetic.bend}{整数の乗算と加算（02\_arithmetic.bend）}
例の関数\code{calculate(x,y)}は\(5x+y\)を計算します。\code{calculate(2,3)}の結果は13です。\code{(x * 5 + y : U32)}では、計算式の後ろに型を示しています。\code{+x}の\code{+}は「値を複数回使える」と示す\term{使用量（quantity）}の指定であり、この短い例では必須ではありません。詳しくは配列とアフィン性の章で扱います。

\section{計算式の型注釈を確認する}
整数加減算の演算子だけを式として書かず、\code{(x + y : U32)}のように型を与える。\code{Nat}と\code{U32}の暗黙変換を当てにしません。\code{U32.to\_nat}・\code{U32.from\_nat}等は\code{bend base}で確認します。\code{U32}には範囲の限界があり、任意精度整数として扱ってはいけません。

\begin{pointbox}
数字だけで動作確認する\code{main() -> U32}は、値の正規化と表示が行われる経路になります。実行可能ファイルを生成する場合は\code{-o}オプションを使います。
\end{pointbox}

\exercise{\code{calculate(0,0)}、\code{calculate(1,2)}、\code{calculate(2,3)}の結果を予想してから実行してみましょう。式の5を10に変えると何が変わるでしょうか。また、定数5を引数\code{factor: U32}に変える方法も考えてみましょう。}

\chapter{パターンマッチで場合分けをする}
\begin{tcolorbox}[colback=blue!3,colframe=blue!16,boxrule=.35pt,arc=1mm,left=3mm,right=3mm,top=2mm,bottom=2mm]
\small Bend 2では、条件に応じて処理を分けるときに\term{パターンマッチ（pattern matching）}を使います。ここでは、真偽値に応じて二つの数字から一つを選ぶ、一般的な例で基本形を学びます。
\end{tcolorbox}

\section{データ型の形を見て処理を選ぶ}
前章で、\code{Status}には\code{Idle\{\}}と\code{Busy\{\}}という値の候補があることを学びました。\term{パターンマッチ（pattern matching）}は、\textbf{値がどの形に当てはまるか}を調べ、形ごとに処理を分ける仕組みです。\code{Bool}の\code{True\{\}}／\code{False\{\}}も同じ考え方で処理できます。
\readcode{13_status_adt.bend}{ユーザー定義型をパターンマッチする例（13\_status\_adt.bend）}
\code{case Idle\{\}:}なら0、\code{case Busy\{\}:}なら1を返します。ここでは数値を返していますが、選択肢の種類が増えたら対応する場合分けも追加する必要があります。

\section{条件に応じて処理を分ける方法}
たとえば「条件が真なら10、偽なら20を返す」という処理には、条件による場合分けが必要です。\code{match}を使うと、それぞれの条件に対応する処理を明確に書けます。

Bend 2の公開ガイドには通常の\code{if ... else}構文がなく、\code{match}と\code{True\{\}}／\code{False\{\}}を使います。特に\textbf{計算した式を直接matchできない}制限があります。たとえば比較式をそのまま\code{match}せず、計算結果を引数として受け取った関数の中で場合分けします。

\readcode{27_match_basics.bend}{真偽値による場合分け（27\_match\_basics.bend）}
\code{choose}に渡す真偽値\code{flag}が\code{True\{\}}なら最初の数値を、\code{False\{\}}なら次の数値を返します。\code{match}は\textbf{引数として受け取った値}に対して行います。条件が何を意味するかは呼び出す側が決めるため、さまざまな処理で再利用できます。

\section{真の場合と偽の場合を追う}
\begin{tabularx}{\linewidth}{@{}c c c X@{}}
\toprule
条件 & 第1候補 & 第2候補 & 返り値\\
\midrule
\code{True\{\}} & 10 & 20 & 10\\
\code{False\{\}} & 10 & 20 & 20\\
\code{True\{\}} & 5 & 9 & 5\\
\bottomrule
\end{tabularx}

\exercise{\code{choose(True\{\},10,20)}と\code{choose(False\{\},10,20)}の返り値を予想して確かめてみましょう。第1候補と第2候補を変えて結果を比較してみましょう。比較式\code{U32.is\_eq}の結果を\code{choose}へ渡す使い方も試してみましょう。}

\chapter{配列の操作とアフィン性}
\begin{tcolorbox}[colback=blue!3,colframe=blue!16,boxrule=.35pt,arc=1mm,left=3mm,right=3mm,top=2mm,bottom=2mm]
\small ここでは\term{配列（array）}を学びます。配列とは、複数の値を順番にまとめ、位置を指定して扱うデータ構造です。あわせて、Bend 2の\term{アフィン性（affinity）}という値の使用方法に関する制約を学びます。
\end{tcolorbox}

\section{配列から値を取り出す基本操作}

配列の要素を読むには、\code{Array.get}を使います。ここで一つだけ、普通の配列操作と違う点があります。
Bend 2では、通常の配列は所有する値として扱われ、操作に渡すと元の変数をそのまま再利用できません。
そこで\code{Array.get}は、\textbf{取り出した要素}だけでなく、\textbf{その後も使うための配列}も組で返します。
まずこの形を見てから、なぜ所有権が必要なのかを確認しましょう。

\code{Array<U32>}は、32ビット符号なし整数を格納する配列型です。\term{種類（kind）}とは「型を分類するための分類」で、\code{Type}は値を自由に複製できない種類、\code{Data}は複製を許せる種類を表します。\code{Array<U32>}は\code{Type}に属するため、同じ配列を二つの処理に勝手に渡せません。読み取るときは\code{Array.get}が\term{組（pair）}として「続けて使う配列」と「読み取った値」の両方を返します。こうすることで同じ配列を勝手に何度も使わず、読み取り後の配列を再び明示して処理できます。

\readcode{04_array.bend}{書き換えと読み出し（04\_array.bend）}
\code{[0 : U32*8n]}は、0を8個並べて作る配列です。\textbf{8要素という長さはこのサンプルが選んだ値}であり、Bend 2の配列が8要素に限定されるわけではありません。\code{Array.set(U32,a,5,42)}は、添字5に42を書き込みます。続く\code{Array.get}は、\code{(a,value)}という組を返します。ただし計算結果をその場で直接分解せず、\code{read\_value}関数の引数\code{r}として受け取り、関数の中で分解しています。最後は\code{value}だけを結果とするので、\code{a}は破棄されます。Bendは\textbf{アフィン}（高々1回の使用）であり、破棄は許されます。後出の「配列操作の糖衣構文」では\code{a[i]}と\code{a[i] <- v}も紹介します。

\begin{cautionbox}
\code{Array}は\code{+a}を宣言して何度も使うことができません。\file{src/12\_affine\_probe\_INVALID.bend}はこの失敗を観察する専用ファイルであり\textbf{正しくないことが正常}。普通の実行対象には含めない。
\end{cautionbox}

\section{値の所有権と利用回数}
\begin{tabularx}{\linewidth}{@{}p{28mm}X@{}}
\toprule
記法 & 解釈\\
\midrule
\code{x: T} & 変数\code{x}を高々1度使う。使わなくてもよい。\\
\code{+x: T} & \code{T}が\code{Data}なら何度も使えます。参照カウント等のコストは無視できません。\\
\code{-A: Data} & 消去される型引数。計算結果の型検査には現れるが実行時には不要。\\
\code{Array<T>} & \code{Type}に属し、通常はコピー可能な\code{Data}ではありません。\\
\bottomrule
\end{tabularx}
\exercise{\code{Array.get}の\code{(array,value)}を\code{(value,array)}と逆にするとどうなるか。\code{Array.set}した配列を使わず、古い配列を再利用すると何が問題か。エラーの全文を記録し、型エラーと規則違反を区別してみましょう。}


\chapter{ListとMaybe：再帰的データと失敗を扱う}
\begin{tcolorbox}[colback=blue!3,colframe=blue!16,boxrule=.35pt,arc=1mm,left=3mm,right=3mm,top=2mm,bottom=2mm]
\small 配列以外のデータの表し方を学びます。\term{リスト（list）}は要素を先頭からつなぐ構造で、\code{Maybe}は「結果があるかどうか」を型に含める仕組みです。
\end{tcolorbox}

\section{リストの構造と配列との違い}
\code{Array}は長さが2の冪などの制約を持つ所有権付き領域であり、\code{List<A>}は\code{Nil\{\}}または\code{Con\{h,t\}}で構成される帰納的な鎖です。処理する数値の並びや、順番に発生するイベントはリストで表現できます。
\readcode{14_list_basics.bend}{リストを構造再帰で処理する（14\_list\_basics.bend）}
\section{MaybeとResultの違い}
\code{Maybe<A>}は値が\code{Some\{value\}}か\code{None\{\}}かを表します。\code{Result}は成功または失敗の情報を保持する用途に適します。文字列を数値に変換する場面では、0を「解析失敗」の印と兼用せず、失敗を型で表す方が安全になります。\code{Maybe}の正確な量パラメータは\code{bend base Maybe}で確認します。
\begin{pointbox}
再帰関数は\term{停止性を証明できる構造的な減少}が重要です。\code{List}では\code{Con}のtail、\code{Nat}では\code{1n+p}のpが自然な減少引数になります。
\end{pointbox}
\exercise{(1) 整数を数個並べた\code{List<U32>}の長さを数えてみましょう。(2) 空リストでは長さが0になることを確かめてみましょう。(3) 数値の解析に失敗した場合を\code{Maybe}で表現する方法を考えてみましょう。}

\chapter{高階関数とクロージャを使う}
\begin{tcolorbox}[colback=blue!3,colframe=blue!16,boxrule=.35pt,arc=1mm,left=3mm,right=3mm,top=2mm,bottom=2mm]
\small 手続き型言語では、関数を「処理を呼び出す手段」と考えることが多いでしょう。ここでは、\textbf{関数そのものを値として受け渡す}考え方を学びます。その際に必要になる\term{高階関数（higher-order function）}と\term{クロージャ（closure）}を区別します。
\end{tcolorbox}

\section{関数も引数にできる}
\term{高階関数}とは、\textbf{関数を引数として受け取るか、結果として関数を返す関数}です。例えば\code{apply}は「どの処理をするか」を引数\code{f}として受け取り、入力\code{x}に適用します。\code{U32 -> U32}は、\code{U32}を受け取り\code{U32}を返す\term{関数型（function type）}です。
\readcode{29_higher_order.bend}{関数を引数に渡す（29\_higher\_order.bend）}
\code{x => (x + 2 : U32)}は\term{無名関数（anonymous function）}です。名前を付けずにその場で関数を作ります。この例では40を受け取り42を返します。手続き型言語の「処理の分岐を増やす」設計とは異なり、\textbf{共通の処理手順を残したまま、実際の計算だけを差し替える}ことができます。

\section{クロージャは外側の値を覚えている}
無名関数が外側の変数を使う場合、関数本体だけでは結果を決められません。そこで、\term{クロージャ}は\textbf{関数と、その関数が利用する外側の値（捕捉した値）をまとめたもの}として動きます。
\readcode{15_closure.bend}{異なる値を捕捉する二つのクロージャ（15\_closure.bend）}
\code{adder(2)}の結果である\code{add2}は、\code{k=2}を覚えた関数です。\code{add2(1)}は3になり、\code{add5(3)}は8になるので、全体では8を返します。\code{adder}は\textbf{関数を返す高階関数}にもなっています。\code{U32.add(2)}のように一部の引数を先に与える\term{部分適用（partial application）}もクロージャの一例です。

\begin{cautionbox}
Bend 2では、\textbf{クロージャはアフィン}です。つまり、同じクロージャを呼べるのは高々1回です。捕捉した\code{U32}が複製可能でも、クロージャ自体を何度も呼ぶことはできません。一方、\code{adder}のようなトップレベルの関数定義は繰り返し呼び出せます。これはPythonやJavaScriptの一般的なクロージャとは重要な違いです。
\end{cautionbox}

\section{一度使う関数と、何度も適用する処理を区別する}
\code{apply(f, x)}は\code{f}を一度使うので、通常の高階関数として書けます。しかし、同じ\code{f}を二度使う\code{f(f(x))}は、一般のアフィンな関数引数では書けません。次章では、\textbf{コンパイル時にコードを埋め込むテンプレート}を学び、この違いを確かめます。
\begin{pointbox}
\term{高階関数}は「関数を受け取る／返す」という\textbf{役割}の名前です。\term{クロージャ}は「関数と捕捉した値」を一体として扱う\textbf{値}の名前です。同じものを指す言葉ではありません。
\end{pointbox}
\exercise{\code{29\_higher\_order.bend}で無名関数を\code{x => (x * 2 : U32)}に変えて結果を予想してみましょう。続いて\code{adder(3)}を作り、捕捉する値によって結果がどう変わるか確認してみましょう。同じクロージャを二度呼ぶ例は別ファイルにして、型検査が拒否する理由も考えてみましょう。}

\chapter{モナドと\code{do}構文を理解する}
\begin{tcolorbox}[colback=blue!3,colframe=blue!16,boxrule=.35pt,arc=1mm,left=3mm,right=3mm,top=2mm,bottom=2mm]
\small この章では、手続き型言語の「一行ずつ処理する」書き方と、関数型言語で\textbf{計算の文脈をつなぐ}書き方の違いを学びます。特に、失敗するかもしれない数値変換を例に\term{モナド（monad）}を理解します。
\end{tcolorbox}

\section{値を返す計算と、失敗し得る計算の違い}
\code{(2 + 3 : U32)}は必ず\code{U32}を返す計算です。一方、文字列を整数に変える\code{U32.read("42")}は成功するとは限りません。\code{"abc"}を渡した場合、返すべき整数がありません。
\term{Maybe}はこの違いを型で示します。成功なら\code{Some\{42\}}、失敗なら\code{None\{\}}です。\code{Maybe<\&2, U32>}という型は、\textbf{U32が得られるかもしれない計算結果}を表します。先頭の\code{\&2}は値の再利用可能性に関するBend 2固有の量指定です。ここでは「値があるかないか」と別の役割であることを押さえてください。

\section{手続き型の失敗判定を\code{Maybe}に置き換える}
従来の言語では、整数を読み取った後に成否を調べ、失敗ならそこで処理を止める\code{if}文を何度も書くことがあります。\code{Maybe}なら\term{失敗の伝播（failure propagation）}を共通の仕組みに任せられます。
\readcode{21_monadic_maybe.bend}{二つの文字列を安全に加算する（21\_monadic\_maybe.bend）}
実行例は\code{add\_strs("40", "2")}が\code{Some\{42\}}、\code{add\_strs("abc", "2")}が\code{None\{\}}です。\textbf{最初の解析に失敗した場合、次の解析と加算へ進みません。}失敗を表すために0や負数などの「特別な数値」を予約する必要もありません。

\section{モナドとは何を共通化した仕組みか}
\term{モナド}は、値を包む型の仕組みと、\textbf{その型を保ったまま複数の計算をつなぐ規則}を組にした抽象化です。Bend 2では、例えば\code{Maybe}に次の二つの操作があります。

\begin{tabularx}{\linewidth}{@{}p{28mm}X@{}}
\toprule 操作 & 役割\\\midrule
\code{M.pure} & 普通の値を\code{M}の計算結果として包みます。\code{Maybe}なら\code{Some\{x\}}に相当します。\\
\code{M.bind} & 包まれた値を次の関数へ渡し、次の計算も\code{M}の形で返します。\code{Maybe}なら\code{None\{\}}を次へ進めず、そのまま返します。\\
\bottomrule
\end{tabularx}

ここで\code{M}は\code{Maybe}や\code{IO}などの型を表す仮の名前です。「モナド＝入出力」でも「モナド＝副作用」でもありません。\code{Maybe}は失敗、\code{Result}は失敗の理由、\code{IO}は外部とのやり取りという\textbf{異なる文脈}を扱います。

\section{\code{bind}を直接書いてみる}
\code{do}は新しい種類の計算を実行する命令ではなく、\code{bind}と\code{pure}を読みやすく書くための\term{構文糖（syntactic sugar）}です。次のコードでは、解析に失敗したときの処理を\code{Maybe.bind}へ任せています。
\readcode{30_monad_explicit.bend}{\code{Maybe.bind}を直接使う（30\_monad\_explicit.bend）}
このプログラムと、前節の\code{do Maybe<\&2, U32>:}を使ったプログラムは同じ結果を返す設計です。\code{x : U32 <- U32.read(a)}は「成功した値を\code{x}へ受け渡して続行する」という意味であり、\textbf{通常の変数代入（\code{x = ...}）とは違います}。

\begin{pointbox}
\code{do M<..., R>:}の\code{x : A <- m}は、\code{M.bind(..., A, R, m, x => ...)}を使う形へ展開されます。\code{return e}は\code{M.pure(..., R, e)}を使う形へ展開されます。ここでの\code{return}は、C言語の\code{return}のような「命令的な早期脱出」と混同しないでください。
\end{pointbox}

\section{モナドの規則と\code{Result}への発展}
一般にモナドとして使う\code{pure}と\code{bind}は、計算を一貫してつなぐために\textbf{左単位律・右単位律・結合律}と呼ばれる規則を満たすように設計します。初めは数式を暗記する必要はありません。「普通の値を包んでから続ける操作」と「そのまま続ける操作」が同じ結果になること、つなぐ場所の括り方を変えても意味が変わらないことが要点です。なお、\textbf{\code{pure}と\code{bind}という名前だけを持つ型が、必ずこの規則を満たすわけではありません}。規則が成り立つかは実装側の責任です。
\code{Result}は\code{Done\{value\}}と\code{Fail\{error\}}を使い、失敗の理由を残せます。\code{Maybe}は単に失敗したことだけを知らせればよいとき、\code{Result}は\textbf{何が間違っていたかを説明したいとき}に向いています。

\section{同じ\code{do}で入出力をつなぐ}
端末への表示では\code{do IO<Unit>:}を使います。\code{IO.print}は「文字列を表示するという\term{効果（effect）}」を記述します。\code{do}内に順番に並べることで、表示の順序を明示できます。\textbf{\code{Maybe}では失敗を、\code{IO}では効果の順序を扱いますが、つなぎ方の構文は共通}です。二つの例で違いを確認してみましょう。
\exercise{\code{add\_strs("12", "30")}と\code{add\_strs("abc", "30")}の結果を予想し、実行してみましょう。\code{do}版と\code{bind}版の結果を比較してみましょう。さらに、三つの文字列を足す関数を作り、途中で失敗した場合の返り値を確かめてみましょう。}

\chapter{テンプレートとコンパイル時の引数}
\begin{tcolorbox}[colback=blue!3,colframe=blue!16,boxrule=.35pt,arc=1mm,left=3mm,right=3mm,top=2mm,bottom=2mm]
\small 「実行するたびに値を受け取る」のと「コンパイル時に処理を埋め込む」のは違います。ここでは\term{テンプレート（template）}を後者として理解します。
\end{tcolorbox}

\section{\code{\textasciitilde{}}は実行時のクロージャではない}
テンプレート・パラメータ\code{\textasciitilde{}f}はコンパイル時にコードとして代入されます。関数適用が複数必要なときでも、アフィンなクロージャを無理にコピーしません。\term{各テンプレート引数の組み合わせごとにインスタンス化}されるため、コードサイズとのトレードオフがあります。
\readcode{16_template.bend}{同じ処理を2度行うテンプレート（16\_template.bend）}
\code{\textasciitilde{}(x => (x + 1 : U32))}は\term{呼び出し元の局所変数を捕捉しない閉じた構文}でなければなりません。テンプレートの宣言順序にも制約があります。\code{List.map}などBaseの定義を\code{bend base}で読むと利用例が分かります。
\section{コンパイル時に処理を選択する}
処理\code{f}を通常の引数にするか、\code{\textasciitilde{}f}としてテンプレート引数にするかは、実行時に選ぶ必要があるか、コンパイル時に決まってよいかで判断します。型と使い回しの制約も比較してみましょう。
\exercise{\code{twice}を\code{thrice}へ拡張してみましょう。テンプレートと普通のクロージャの寿命・生成コードの違いを、実行時とコンパイル時に分けて説明してみましょう。}

\chapter{数値・文字列・演算子と標準ライブラリ}
\begin{tcolorbox}[colback=blue!3,colframe=blue!16,boxrule=.35pt,arc=1mm,left=3mm,right=3mm,top=2mm,bottom=2mm]
\small 数値型と文字列型は、用途ごとに使い分けます。まず整数と小数の違い、演算子が受け取る型、標準ライブラリの調べ方を整理します。
\end{tcolorbox}

\section{主要な値の型}
Bend 2の基本的な数値型は\code{Nat}、\code{U32}、\code{F32}です。\code{F32}は小数を近似的に表す\term{浮動小数点数（floating-point number）}で、ビット単位の演算や丸めに注意します。\code{Bool}、\code{Cmp}、\code{Char}、\code{String}も使用します。現行の公式ガイドは\term{F64やI64が最初から使える}とは述べていない。
\readcode{17_numbers_strings.bend}{整数演算と文字列化（17\_numbers\_strings.bend）}
\section{演算子の型付けとビット演算}
\code{(a + b : U32)}のように、型注釈は\term{その演算を囲う式自身}に置きます。\code{\&\&}・\code{||}はBool、\code{++}はString。\code{.\&.}、\code{.|.}、\code{.\^{}.}や\code{<<}などは\code{U32}のビット演算で、一般の数学的な等式と区別します。
\section{標準ライブラリを調べる}
\code{import Base}で読み込む\term{標準ライブラリ（standard library）}には、文字列化や配列操作など、よく使う関数が用意されています。\code{bend base U32}で整数の操作を、\code{bend base String}で文字列の操作を調べられます。例えば、整数を文字列で表示したいときは\code{U32.show}を使います。実行して確認してみましょう。
\begin{lstlisting}[caption={標準ライブラリの確認例},language={}]
bend base U32
bend base String
bend src/17_numbers_strings.bend
\end{lstlisting}
\exercise{(1) \code{Nat}と\code{U32}が自動変換されない例を探してみましょう。(2) \code{U32}のAND・OR・XORと比較の違いを小さい入力で表にまとめてみましょう。(3) 数値を文字列へ変換するときに使用できる関数を調べてみましょう。}

\chapter{モジュールと複数ファイルのプログラム}
\begin{tcolorbox}[colback=blue!3,colframe=blue!16,boxrule=.35pt,arc=1mm,left=3mm,right=3mm,top=2mm,bottom=2mm]
\small 一つのファイルが長くなったら、機能ごとに分割します。\term{モジュール（module）}とは、関数や型をひとまとめにした再利用単位です。
\end{tcolorbox}

\section{\code{import ./... as M}}
ソースを複数のファイルに分けると、数値計算、入出力、テストなどの責務を整理しやすくなります。各\code{.bend}ファイルがモジュールです。\code{import ./file.bend as M}は、取り込んだ定義を\code{M}という名前で区別できるようにします。以下の例では\code{module\_math.bend}を取り込みます。
\readcode{module_math.bend}{取り込まれる側の定義（module\_math.bend）}
\readcode{19_module_main.bend}{別ファイルから定義を呼ぶ（19\_module\_main.bend）}
\begin{lstlisting}[language={},caption={モジュールの実行と将来の分割}]
bend src/19_module_main.bend
# 例: math.bend, parser.bend, main.bend
\end{lstlisting}
\section{モジュール名を使って関数を呼ぶ}
同じ名前の関数が別のモジュールにある場合でも、\code{M.square}のようにモジュール名を付けて呼び出せます。
\exercise{\code{square}だけを別の\code{.bend}へ切り出し、\code{import}を通して呼べることを確認してみましょう。また、ファイル同士が相互に取り込まない構成を考えてみましょう。}

\clearpage
\part{ライフゲームを設計して組み立てる}
\partchaptertoc
\begin{pointbox}
最初にライフゲームの規則を確認し、盤面の配列、初期配置、近傍の数え方、世代更新、表示、テストを実装します。Part Iで学んだ小さな処理を組み合わせて完成させましょう。
\end{pointbox}
\chapter{ライフゲームのしくみと制作するもの}\label{ch:life-spec}
\begin{tcolorbox}[colback=blue!3,colframe=blue!16,boxrule=.35pt,arc=1mm,left=3mm,right=3mm,top=2mm,bottom=2mm]
\small 本章では、ライフゲームの規則を日本語と数式で整理します。\term{仕様（specification）}とは、プログラムが満たすべき動作の条件です。先に条件を決めておくと、テストで確認できます。
\end{tcolorbox}

\section{何を作るか}
ライフゲームは二次元格子の各セルが生存（1）か死亡（0）かを持つ\term{セル・オートマトン（cellular automaton）}です。周囲8方向にある\term{近傍セル（neighboring cell）}の生存数を数えて、次の状態を決めます。簡単な規則ですが、移動する模様や形を変え続ける模様が生まれます。そのため、小さなプログラムでも\term{モデル化（modeling）}、境界での処理、テスト、性能評価を学べます。

本書では格子を8行×8列、セルを\code{U32}の0と1、座標を\((x,y)\)（左上が\((0,0)\)）とします。盤面の右端は左端に、下端は上端につながります。これを\term{周期境界条件（periodic boundary condition）}と呼びます。全セルを\textbf{同時に更新}します。

\section{生死規則 B3/S23}
セル\(a\in\{0,1\}\)と、その8近傍の生存数\(n\in\{0,\ldots,8\}\)に対して
\[
F(a,n)=\begin{cases}
1&n=3\quad\text{（生死によらず誕生／生存）}\\
a&n=2\quad\text{（現状維持）}\\0&\text{その他}
\end{cases}
\]
これが\textbf{B3/S23}（Birth at 3, Survive at 2 or 3）です。いわゆる「生存セルは2～3近傍で生き残る」という説明と矛盾しません。生存数だけでは次状態は決まらず、\(n=2\)なら現在の状態を参照することに注意してください。

\begin{center}
\begin{tikzpicture}[scale=.62]
\foreach \x in {0,...,7}{\foreach \y in {0,...,7}{\draw[gray!50] (\x,\y) rectangle ++(1,1);}}
\foreach \x/\y in {1/7,2/6,0/5,1/5,2/5}{\fill[teal!75] (\x+.07,\y+.07) rectangle ++(.86,.86);}
\node[anchor=west,align=left,font=\sffamily\small] at (8.5,5.7) {世代0のグライダー\\生存: 5セル\\1=生存、0=死亡};
\draw[-{Stealth},very thick,teal] (8.4,4.3)--(10.0,3.4);
\node[anchor=west,font=\sffamily\small] at (8.6,3.1) {4世代後に右下へ1マス};
\end{tikzpicture}
\end{center}

\section{受け入れ条件}
\begin{enumerate}
\item 正しいグライダーを初期状態として5枚（世代0～4）の8×8盤面を表示します。
\item 近傍として数えるのは周囲の8マス\textbf{だけ}です。中心のセル自身は数えません。
\item すべてのセルの新状態を、更新前の盤面だけから計算します。\term{その場更新（in-place update）}を1セルずつ行う方法は採用しません。
\item ブロックが静止、ブリンカーが周期2、グライダーが4世代で斜めへ1マス動くことをPythonの独立実装と照合します。
\item 表示と更新を分離して、後で表示形式・実行先を交換できる形にします。
\end{enumerate}
\exercise{上の規則から、現在0/1の各場合について近傍数0～8の真理値表（\(2\times9\)個）を作ってみましょう。なぜ「3なら1、2なら現状、それ以外0」の3分岐で全部を処理できるのかを説明してみましょう。}

\chapter{一次元配列で8×8盤面を表す}
\begin{tcolorbox}[colback=blue!3,colframe=blue!16,boxrule=.35pt,arc=1mm,left=3mm,right=3mm,top=2mm,bottom=2mm]
\small 前章で決めたライフゲームの規則を関数にし、続いて二次元の盤面を一次元の配列に並べます。ここから実際のシミュレータを組み立てます。
\end{tcolorbox}

\section{生死規則を関数で表す}

\readcode{03_rule.bend}{ライフゲームの生死規則（03\_rule.bend）}
\code{rule(alive,neighbors)}では、\code{neighbors}が3なら1、2なら\code{alive}、それ以外なら0を返します。\code{choose}の内部では、Part Iで学んだ\code{match}を使っています。\code{+neighbors}は同じ数値を複数回比較するための使用量指定です。
\exercise{\code{rule(0,3)}、\code{rule(1,2)}、\code{rule(1,4)}の返り値を確認してみましょう。現在の状態0/1と近傍数0～8の全18通りを表にして、期待値と一致するかを調べてみましょう。}

\section{二次元の盤面を一次元配列に並べる}
ボードを二重配列にするのではなく、長さ64の一次元\code{Array<U32>}とします。変換は\(i=8y+x\)、逆変換は\(x=i\bmod8,\ y=\lfloor i/8\rfloor\)。範囲は\(0\leq x,y<8\)、\(0\leq i<64\)です。配列アクセスに用いる\code{U32}添字は固定幅として扱います。


\begin{pointbox}
8行×8列を一次元で表現するため、ここでは64要素の配列を作ります。
Bend 2 の基本的な配列は要素数が2の累乗になる平衡二分木で表されます。
たとえば\code{[0 : U32*256n]}は256要素、\code{Array.new(U32, 8n, 0)}も深さ8の256要素です。
\code{Array.new}の深さは要素数そのものではなく、要素数\(2^d\)の指数\(d\)です。
また\code{Array<Array<U32>>}のように\textbf{配列を要素に持つ配列}で二次元構造を表せます。
ただし、本書では座標計算と更新を見通しよくするために一次元に並べます。
縦横の長さが2の累乗でない盤面は、余白を設けたり、別の構造を選んだりする設計が必要です。
\end{pointbox}

\section{すべてのセルを0で初期化する}
\code{[0 : U32*64n]}で要素64個を確保します。書き込みは\code{Array.set(U32,b,index,value)}で新しい所有者に移る。よって\code{b0,b1,b2,...}と変数名を変えていくとアフィン性が見えます。インデックスは\(8y+x\)で求めます。

\begin{center}
\begin{tikzpicture}[x=.78cm,y=.62cm]
\foreach \x in {0,...,7}{\foreach \y in {0,...,3}{
\draw[soft,thick] (\x,\y) rectangle ++(1,1);
\pgfmathtruncatemacro{\idx}{8*(3-\y)+\x}
\node[font=\ttfamily\tiny] at (\x+.5,\y+.5) {\idx};
}}
\node[anchor=west,font=\sffamily\small] at (8.4,1.5) {添字：上段が0～7、次が8～15…};
\end{tikzpicture}
\end{center}
\section{座標を配列の添字に変換する}
Bendの配列は範囲外インデックスがラップする仕様を持ちます。しかし\textbf{一次元のラップ}は\textbf{二次元のトーラス隣接}と同じではありません。\((7,0)\)の右隣は\((0,0)\)であって\((0,1)\)ではありません。近傍計算は座標ごとに\(x\)と\(y\)を循環させてから一次元化します。

\exercise{添字7と8の意味を説明してみましょう。\((7,0)\)から右に動いたとき\code{i+1}してはいけない理由を図にして考えてみましょう。}

\chapter{グライダーを初期配置する}
\begin{tcolorbox}[colback=blue!3,colframe=blue!16,boxrule=.35pt,arc=1mm,left=3mm,right=3mm,top=2mm,bottom=2mm]
\small 次に、最初の盤面を用意します。小さな生存セルの模様を指定し、正しい位置に置けたかを自分で確認します。
\end{tcolorbox}

\section{小さなパターンを正確に配置する}
グライダーは生存セル5個で右下へ移動する周期4の模様。本書の初期配置は次の5点です。\[\{(1,0),(2,1),(0,2),(1,2),(2,2)\}\]行優先添字に変換すると\(1,10,16,17,18\)になります。

\readcode{05_glider.bend}{初期グライダー（05\_glider.bend）}
図で見れば次の配置です。\code{1}を生存、\code{0}を死亡とします。
\begin{lstlisting}[language={},numbers=none,caption={世代0の配置（8列×8行）}]
01000000
00100000
11100000
00000000
00000000
00000000
00000000
00000000
\end{lstlisting}
\section{盤面を数える}
グライダーの生存数は5。今はセル値を\code{U32}として持つため、\code{2}や\code{42}も表現自体はできます。\textbf{セルが0/1であるという不変条件}は今のところシードと更新関数がそれ以外を作らないという開発規約に依存しています。後半では「\code{U32}で表現可能」と「型で0/1に限定される」を分けて考えます。

\exercise{ブロック（\((3,3),(4,3),(3,4),(4,4)\)）を初期化する\code{block()}を作ってみましょう。ブリンカーを\(y=3,x=3,4,5\)に置く\code{blinker()}を作ってみましょう。出力で指定した場所だけが1となることを確認してみましょう。}

\chapter{トーラスの座標と8近傍}
\begin{tcolorbox}[colback=blue!3,colframe=blue!16,boxrule=.35pt,arc=1mm,left=3mm,right=3mm,top=2mm,bottom=2mm]
\small どのセルが隣かを決めます。今回は画面の右端と左端、上端と下端がつながる\term{トーラス（torus）}状の盤面です。
\end{tcolorbox}

\section{盤面の端を越えた座標を計算する}
左隣は\((x-1+8)\bmod8\)、右隣は\((x+1)\bmod8\)。\(x\in\{0,\ldots,7\}\)に限れば\((x+7)\bmod8\)で代用できます。同様に上下も\((y+7)\bmod8\)、\((y+1)\bmod8\)で求まる。負の数を\code{U32}に入れない。
\[
L(x)=(x+7)\bmod8,\quad R(x)=(x+1)\bmod8,\\
T(y)=(y+7)\bmod8,\quad B(y)=(y+1)\bmod8.
\]
\section{周囲8セルを読み取り生存数を数える}
\readcode{06_neighbors.bend}{近傍数を足す（06\_neighbors.bend）}
\code{neighbors}は\code{(board,sum)}を返します。元の配列を\code{b1,b2,...,b8}へ順に渡すので、8回の読み取りで\code{board}を違法に共有しません。\code{+xl}、\code{+xr}、\code{+yt}、\code{+yb}は、計算した座標値を複数の読み取りに使うために必要。

\begin{pointbox}
近傍の重複を防ぐには、\(3\times3\)の中心\((x,y)\)を除く8箇所を列挙し、同じ\((x,y)\)を含めていないか紙に書いて点検します。\textbf{旧盤面から読み、読んだ配列を戻す}という型付きの設計は、この過程の所有権をはっきりさせる。
\end{pointbox}

\section{境界にいるセルを手計算する}
例えば\((0,0)\)の8近傍は\((7,7),(0,7),(1,7),(7,0),(1,0),(7,1),(0,1),(1,1)\)。四隅のセルで座標の循環をテストします。\((4,4)\)など内部座標だけを調べても境界バグは見つからない。

\exercise{(1) \((0,0)\)の8近傍座標を書き出してみましょう。(2) 初期グライダーに対して\((1,1)\)の生存近傍数を手計算してみましょう。(3) コード内で\code{+xl}を外して、どの使用が重複するか確認してみましょう。(4) \file{06\_neighbors.bend}で\code{neighbors(4,4,...)}の期待値は3になることを考えてみましょう。}

\chapter{全セルを同期的に更新する}
\begin{tcolorbox}[colback=blue!3,colframe=blue!16,boxrule=.35pt,arc=1mm,left=3mm,right=3mm,top=2mm,bottom=2mm]
\small いよいよ1世代の更新を実装します。\term{同期更新（synchronous update）}とは、次の盤面を計算するとき、すべてのセルが同じ古い盤面を読む方法です。
\end{tcolorbox}

\section{更新に必ず二枚の盤面を使う}
\term{正しい更新}では、時刻\(t\)の状態\(B_t\)だけを見て\(B_{t+1}\)を作る。例えば上から順に古い配列に直接書き込むと、後で処理するセルが更新済みの近傍値を読みます。その瞬間、同時更新ではなく走査順依存の別のゲームになります。
\[
B_{t+1}(x,y)=F\Big(B_t(x,y),\sum_{(u,v)\in N(x,y)}B_t(u,v)\Big).
\]
\section{再帰処理が終わる条件を決める}
\term{燃料（fuel）}とは、再帰を続ける残り回数を表す値を比喩的に呼んだものです。燃料が0になれば止まる、と決めれば無限に繰り返さないことを説明しやすくなります。

\code{fill(todo, i, old, out)}の\code{todo}は\code{Nat}で、残り何セル更新するかを表します。\code{match todo}で0なら出力盤面を返し、\code{1n+p}なら1セル更新して\code{fill(p,...)}を呼びます。再帰が必ず止まることを\code{todo}の減少から説明できます。

\readcode{07_life_core.bend}{完成した純粋な更新核（07\_life\_core.bend）}

\section{ソースを読む順番}
\begin{enumerate}
\item \code{glider()}は5セルを1にした入力の配列を作る。
\item \code{step(board)}は旧盤面と全0の新盤面を用意します。
\item \code{fill(64n,0,old,out)}はセル\(i=0\)から\(63\)を処理します。
\item \code{neighbors}は旧盤面を読み、\code{Array.get}は対象セルの現在値を読みます。
\item \code{rule}で次状態を決め、\code{Array.set}で新盤面の対応位置に保存します。
\item \code{fill}が残りを処理します。旧盤面は最後まで変更しません。
\end{enumerate}
この走査は計算量\(O(WH)\)、追加記憶量\(O(WH)\)です。8方向参照の定数が大きいため、厳密な性能比較では「セル1個当たりの読み込み・格納」を数える。

\section{別の実装と結果を比較する}
同じ\code{neighbors}を使ってテスト期待値を作ると、その関数に同じ誤りがあった場合に見逃します。したがって\file{scripts/reference.py}は、別の記述・別の実装言語で定義した生存セルの\code{set}から計算します。別の方法で得た期待値と照合すると、見落としていた誤りを発見しやすくなります。

\exercise{(1) \code{fill}の先頭引数を\code{U32}のカウンタに置き換えると停止性検査にどんな障害があるか。(2) \code{step}を旧配列1枚への更新に変え、グライダーの軌跡がどこから破綻するか観察してみましょう。(3) \code{step(glider())}の8×8値をオラクルと照合してみましょう。}

\chapter{複数世代の実行とターミナル描画}
\begin{tcolorbox}[colback=blue!3,colframe=blue!16,boxrule=.35pt,arc=1mm,left=3mm,right=3mm,top=2mm,bottom=2mm]
\small 1世代だけでなく複数世代を動かし、ターミナルへ表示します。計算と表示を別々の関数にすると、後から見た目だけを変更できます。
\end{tcolorbox}

\section{世代を繰り返す関数}
\code{evolve(steps,board)}は、0ステップならそのまま返し、正の\code{Nat}なら\code{step}を呼んでから\code{evolve(p,...)}を呼びます。ここで\code{p}が1減った残りの自然数になっています。型検査は停止の根拠を要求するため、\code{evolve}の再帰が無限にループする例は通常許されません。

\section{出力を作る関数は更新と分ける}
一つの行は8セルを読み文字列化し、8行を結合します。\code{render\_cells}と\code{render\_rows}は\code{(board,String)}を返します。読み出しに使った配列を次の処理へ回せるようにするためです。\code{IO.print}は\code{show}だけが持つ。

\readcode{08_life_terminal.bend}{完成版：5世代を表示するシミュレータ（08\_life\_terminal.bend）}

\section{実行する}
\begin{lstlisting}[language={},caption={完成版の実行}]
bend src/08_life_terminal.bend --check-only
bend src/08_life_terminal.bend
bend src/08_life_terminal.bend -o build/life
./build/life --threads 1
\end{lstlisting}
このプログラムは\term{世代0～4の5枚}を表示します。\code{IO.print}が末尾の改行を追加するので、盤面の間には空行が入る。\code{show(5n,0,glider())}を\code{show(10n,0,glider())}にすれば10枚になります。

\section{期待される世代4}
\begin{lstlisting}[language={},numbers=none,caption={Python参照実装が算出した世代4の一部}]
00100000
00010000
01110000
\end{lstlisting}
これはグライダーが右下へ1マス移動した世代4の盤面から3行を抜き出したものです。完全な8行は\file{scripts/reference.py}で確認できます。

\exercise{(1) \code{show}の世代数を3、5、10に変更してみましょう。(2) 表示を\code{1/0}から\code{\#/\.}へ変えるための\code{glyph}関数を追加してみましょう。(3) 初期配置の章で作った\code{block()}を使い、静止パターンであることを調べてみましょう。}

\chapter{テスト設計と失敗の切り分け}
\begin{pointbox}CI（継続的インテグレーション）は、GitHubへの変更をきっかけにテストを自動実行する仕組みです。テストは計算結果の不一致を検出し、失敗したときには終了コードを0以外にして通知します。\end{pointbox}

\begin{tcolorbox}[colback=blue!3,colframe=blue!16,boxrule=.35pt,arc=1mm,left=3mm,right=3mm,top=2mm,bottom=2mm]
\small 動作を確認するための\term{テスト（test）}を作ります。完成したプログラムと、別方式で作った参照実装の結果を比べます。
\end{tcolorbox}

\section{プログラムの正しさをテストで確認する}
ライフゲームの出力が期待どおりか確認しましょう。空の盤面や静止する模様など、結果が分かっている入力を用意して比べます。

コンパイラが型の整合性と停止性を確認しても、\code{neighbors}で中央セルを足してしまう誤りや、\code{rule}で2と3を取り違える誤りは型だけでは防げない。少なくとも\term{局所規則}、\term{初期配置}、\term{時刻発展}、\term{境界条件}の4レベルで検査します。

\section{参照実装を動かす}
\begin{lstlisting}[language={},caption={付属のPython独立テスト（標準ライブラリのみ）}]
python3 -m unittest discover -s tests -v
python3 scripts/reference.py
python3 scripts/check_bend.py
python3 scripts/check_bend.py --native
\end{lstlisting}
最後の2行はBendが導入されている場合の任意の\term{統合テスト}です。\code{check\_bend.py}はBendの表示から5枚の盤面を取り出し、Pythonオラクルの出力と改行込みで比較します。失敗したらdiffを表示します。

\section{用意された参照テスト6件}
\begin{tabularx}{\linewidth}{@{}p{49mm}X@{}}
\toprule
テスト & 検査内容\\
\midrule
\code{test\_empty\_stays\_empty} & 空盤面は空のまま。\\
\code{test\_block\_is\_still\_life} & ブロックは1世代後も変化しません。\\
\code{test\_blinker\_oscillates} & ブリンカーは横 $\rightarrow$ 縦 $\rightarrow$ 横の2周期。\\
\code{test\_glider\_moves...} & 4世代後、グライダーが斜めに1進みます。\\
\code{test\_alive\_count...} & グライダーの生存数は12世代を通じて5。\\
\code{test\_transcript...} & 出力の枚数・終端テキストの形。\\
\bottomrule
\end{tabularx}

\begin{pointbox}
Pythonの参照テスト6件と、Bend 2.0.37での5世代全セルの一致確認はCIで成功しています。コンパイラを更新した場合は、テストを再実行しましょう。
\end{pointbox}

\exercise{(1) トーラス端をまたぐブリンカーを設定し、2周期を検証するテストを書いてみましょう。(2) \code{rule}の「2のとき現状維持」を意図的に崩し、どのテストが落ちるかを報告してみましょう。(3) 手計算・Python・Bendの3列からなる検証表を作ってみましょう。}


\chapter{独立した計算をCPUで並列に実行する}
\begin{tcolorbox}[colback=blue!3,colframe=blue!16,boxrule=.35pt,arc=1mm,left=3mm,right=3mm,top=2mm,bottom=2mm]
\small これまでの\code{step}は1世代を順番に更新しました。ここでは二つの独立した計算を同時に進める\term{並列計算（parallel computation）}を学びます。
\end{tcolorbox}

\section{二つの関数呼び出しを同時に始める}
手続き型言語では、一つの関数の結果を待ってから次の関数を呼ぶ書き方が一般的です。Bend 2では、互いに独立した呼び出しを次のように横に並べます。
\begin{lstlisting}[caption={二つの独立した呼び出し}]
a b = f(x) g(y)
# a と b の両方が得られた後で次の処理へ進みます。
\end{lstlisting}
これを\term{並列呼び出し（parallel call）}と呼びます。\textbf{二つの処理が同じ値を不正に共有しないこと}はBendの純粋性とアフィン性に支えられます。仕事量が大きく違うと、片方が終わるまで待つ時間が増えるので、負荷の均等さも重要です。

\readcode{09_parallel_tree.bend}{二つに分割しながら値を足す（09\_parallel\_tree.bend）}
\code{pow2(12n)}の結果は\(2^{12}=4096\)です。\code{match}で再帰が終了する条件を定め、二つの再帰呼び出しを同じ深さで実行しています。

\section{ライフゲームを二つ同時に計算する}
ライフゲームには、\textbf{別々の初期盤面から始まる二つのシミュレーション}という独立した仕事があります。異なる入力を使うので、所有権を共有しなくても並列に実行できます。
\readcode{31_parallel_life_pair.bend}{独立した二つの盤面を並列に更新する（31\_parallel\_life\_pair.bend）}
この例は\textbf{複数盤面の同時計算}です。一つの盤面の64セルすべてを並列化しているわけではありません。\code{Array<U32>}を再利用可能な値として二つの処理に同時に渡すことはできないので、単一盤面内を並列化する場合には入力の分割や複製を工夫する必要があります。

\begin{lstlisting}[language={},caption={ネイティブCPU実行でスレッド数を比較する}]
bend src/09_parallel_tree.bend -o build/parallel_tree
./build/parallel_tree --threads 1
./build/parallel_tree --threads 4
\end{lstlisting}
\term{スレッド（thread）}とは、CPU上で仕事を実行する単位です。結果が同じことを確認した後で、所要時間を測りましょう。小さな入力では並列化の管理にかかる時間が勝つこともあります。
\exercise{\code{pow2(10n)}と\code{pow2(12n)}を実行して、結果を予想してみましょう。二つの独立した初期盤面について、並列計算した結果と逐次計算した結果が一致するか確かめましょう。}

\chapter{GPUで計算を実行する}
\begin{tcolorbox}[colback=blue!3,colframe=blue!16,boxrule=.35pt,arc=1mm,left=3mm,right=3mm,top=2mm,bottom=2mm]
\small \term{GPU（graphics processing unit）}は多数の計算を同時に実行できる装置です。Bend 2では、並列な関数呼び出しの計算先を指定するための短い記法があります。
\end{tcolorbox}

\section{関数名の後ろの\code{!}を読む}
\begin{lstlisting}[caption={通常の呼び出しとGPU呼び出し}]
value_cpu = pow2(12n)
value_gpu = pow2!(12n)
\end{lstlisting}
\code{!}は論理否定ではなく、\term{GPU呼び出し（GPU call）}を表します。GPUが利用できない環境ではCPU側で実行されることがあります。JavaScript実行環境ではこの指定が逐次実行として扱われる場合があります。\textbf{GPU呼び出しを書くだけで必ず速くなるわけではありません。}

\readcode{10_gpu_tree.bend}{GPU実行を要求する例（10\_gpu\_tree.bend）}
\readcode{32_gpu_life_call.bend}{ライフゲームの計算をGPUへ渡す書き方（32\_gpu\_life\_call.bend）}
後者の\code{evolve!}は\textbf{GPUへの実行指定を学ぶための例}です。その内部の\code{step}は逐次処理なので、セル単位のGPU並列版が完成したことにはなりません。単一盤面を真に並列更新するには、旧盤面の各セルを安全に読むデータ分割と、計算後の盤面を統合する設計が必要です。

\section{GPUの有無を分けて確かめる}
\begin{lstlisting}[language={},caption={GPU実行とCPU代替実行}]
bend src/10_gpu_tree.bend -o build/gpu_tree
./build/gpu_tree
./build/gpu_tree --gpu off
\end{lstlisting}
ネイティブ実行に必要なGPUドライバや対応コンパイラは環境によって異なります。まず出力が4096で一致するか確かめ、その後で速度を比べましょう。8×8の小さな盤面では、GPUの起動コストが大きくなる可能性があります。
\exercise{CPU実行とGPU指定の違いを説明してみましょう。\code{!}を付けたプログラムについて、計算結果が同じかどうか調べてみましょう。速度を測った場合は、使用したハードウェアと測定方法も記録しましょう。}

\chapter{ライフゲームの規則を形式的に証明する}
\begin{tcolorbox}[colback=blue!3,colframe=blue!16,boxrule=.35pt,arc=1mm,left=3mm,right=3mm,top=2mm,bottom=2mm]
\small テストは選んだ入力を調べます。一方、\term{形式的証明（formal proof）}は「すべての対象入力について成り立つ」という主張を、検査できる形で示します。Bend 2では\code{law}と\code{def}でそれを記述できます。
\end{tcolorbox}

\section{法則と証明を対応させる}
例えば「近傍が3つなら必ず次の状態は生存」という\term{命題（proposition）}を考えます。
\begin{lstlisting}[caption={法則と証明を記述する最小例}]
import Base
law three_is_alive:
  for cell: U32
  {rule(cell, 3) == 1 : U32}

def three_is_alive(cell):
  {==}
\end{lstlisting}
上の例は\code{rule}の定義を省略しています。実際に実行する際は次のファイルを使用してください。
\readcode{11_rule_proofs.bend}{生死規則に対する三つの形式的証明（11\_rule\_proofs.bend）}
\code{for cell: U32}は、対象を\textbf{すべての\code{U32}の値}に広げます。\code{\{a == b : U32\}}は二つが等しいという\term{等式の型（equality type）}で、\code{\{==\}}は左右が計算により同じ式になることを示す証明です。

このファイルにある三つの法則の意味は、近傍数が3なら出力が1、2なら元の状態を保つ、0なら出力が0、です。これらは固定した18例より広い条件について証明しますが、\textbf{B3/S23のすべての分岐が証明済みという意味ではありません}。

\section{帰納法と証明の検査}
\term{数学的帰納法（mathematical induction）}は、最初の値で成り立つことと、一つ小さい値で成り立つなら次の値でも成り立つことを示す方法です。Bendでは\code{match}と再帰を使って書けます。
\readcode{24_proof_induction.bend}{自然数を用いた帰納法と書き換え（24\_proof\_induction.bend）}
\code{\%plus\_zero(p)}は、再帰で得た証明を用いた\term{等式の書き換え（rewriting）}です。\code{?TODO}のような未完成の箇所を残したものを証明完了とは扱いません。
\begin{lstlisting}[language={},caption={Bendの証明検査}]
bend src/11_rule_proofs.bend --verdict
bend src/24_proof_induction.bend --verdict
\end{lstlisting}

\section{盤面全体の仕様を証明するには}
一つのセルの\code{rule}を証明しても、\code{neighbors}がちょうど8方向を数え、\code{step}が旧盤面だけを使って64セルを正しく更新したことまでは分かりません。全体を証明するには、\textbf{添字変換、近傍の列挙、全セルの同期更新}について別の法則を作り、組み合わせる必要があります。本書では、局所規則の形式的証明と、盤面全体の実行テストを区別します。

\exercise{近傍数が1のとき必ず死ぬことを表す\code{law one\_dies}を追加し、\code{\{==\}}で証明してみましょう。次に、近傍を数える関数の正しさを法則として書くには、どの情報が必要か考えてみましょう。}

\chapter{完成版の設計をレビューする}
\begin{tcolorbox}[colback=blue!3,colframe=blue!16,boxrule=.35pt,arc=1mm,left=3mm,right=3mm,top=2mm,bottom=2mm]
\small いったん完成したプログラムを読み直します。\term{保守性（maintainability）}とは、後から機能追加や修正をしやすいことです。
\end{tcolorbox}

\section{関数が担う責務}
\begin{center}
\begin{tikzpicture}[>=Stealth,node distance=10mm, every node/.style={font=\sffamily\small}]
\node[draw=blue,fill=blue!5,rounded corners] (seed) {glider: 初期配置};
\node[draw=blue,fill=blue!5,rounded corners,right=of seed] (step) {step: 1世代更新};
\node[draw=teal,fill=teal!7,rounded corners,right=of step] (view) {render: 文字列化};
\node[draw=warm,fill=warm!7,rounded corners,below=of step] (show) {show: 繰返しとIO};
\draw[->] (seed)--(step);\draw[->] (step)--(view);
\draw[->] (step)--(show);\draw[->] (view)--(show);
\node[draw=ink,fill=paper,rounded corners,above=of step] (rule) {neighbors + rule: 局所計算};
\draw[->] (rule)--(step);
\end{tikzpicture}
\end{center}
\code{rule}は1セルの遷移、\code{neighbors}は8セルの集計、\code{step}は盤面全体の遷移、\code{render}は表示、\code{show}は実行制御を担当します。それぞれの関数の役割が明確だと、テストで誤りの場所を絞り込めます。

\section{保守性チェック}
\begin{enumerate}
\item \textbf{定数が散らばっている}: 8と64が複数箇所にあります。定数に名前を付けると変更箇所を減らせます。
\item \textbf{0/1の意味が型に表れていない}: 新しく生死状態を表す型を定義する方法もあります。
\item \textbf{表示が計算から独立}: 文字列化を別の関数にしているので、セルの表現を変更しても世代更新の関数を修正せずに済みます。
\item \textbf{出力枚数は固定}: 表示する世代数は\code{show}への引数で変更できます。

\end{enumerate}
\exercise{\code{glider}、\code{step}、\code{render}、\code{show}について、引数・返り値・画面表示を行うかどうかを表にまとめてみましょう。\code{rule}の実装を変更したとき、どのテストを実行すればよいでしょうか。}


\clearpage
\part{演習の解答と実行リファレンス}
\partchaptertoc
\chapter{重要演習の解答と検証方法}
\begin{tcolorbox}[colback=blue!3,colframe=blue!16,boxrule=.35pt,arc=1mm,left=3mm,right=3mm,top=2mm,bottom=2mm]
\small ここからは解答・付録です。自力で解いてから照合し、単に答えが一致したかだけでなく、理由も説明できるか確認してください。
\end{tcolorbox}

\section{18ケースの局所規則}
\begin{tabular*}{\linewidth}{@{\extracolsep{\fill}}c ccccccccc@{}}
\toprule
\(n\) & 0&1&2&3&4&5&6&7&8\\
\midrule
\(F(0,n)\)&0&0&0&1&0&0&0&0&0\\
\(F(1,n)\)&0&0&1&1&0&0&0&0&0\\
\bottomrule
\end{tabular*}
\section{代表的な添字}
\((0,0)\mapsto0,\ (7,0)\mapsto7,\ (0,7)\mapsto56,\ (7,7)\mapsto63\)。\((2,3)\mapsto26\)。\((7,0)\)の右は\((0,0)\)で、\((0,1)\)ではありません。
\section{グライダーの変位}
世代0の集合は\(\{(1,0),(2,1),(0,2),(1,2),(2,2)\}\)。4回の更新後は\(\{(2,1),(3,2),(1,3),(2,3),(3,3)\}\)。\code{tests/test\_reference.py}にはこの性質を全セル集合として比較するテストがあります。
\section{近傍読み出し回数}
1世代で\(N=WH\)個のセルを更新するなら、現行の各セルは\code{neighbors}の8回と現在セルの1回で、合計\(9N\)回の\code{Array.get}を使います。\code{Array.set}は\(N\)回。計算量の章の「8方向参照」と現在セルの読取りを区別すること。
\section{典型的なバグ}
\begin{tabularx}{\linewidth}{@{}L{55mm}X@{}}
\toprule
症状 & 疑う箇所\\
\midrule
盤面右端付近だけ違う & \(x\)方向のラップと一次元添字の混同。\\
ブリンカーが1周期で崩れる & 旧世代を更新しながら後続の近傍を読んでいる。\\
すべてのセルが1になる & 自分自身を近傍として含めた、あるいは条件式の3と2を間違えた。\\
\code{used more than once} & 配列を戻り値として引き継いでいない、または\code{+}の不足。\\


\bottomrule
\end{tabularx}

\chapter{環境・コマンド・参照先}
\begin{tcolorbox}[colback=blue!3,colframe=blue!16,boxrule=.35pt,arc=1mm,left=3mm,right=3mm,top=2mm,bottom=2mm]
\small 困ったときに参照するコマンドをまとめています。\term{再現性（reproducibility）}は、同じ条件で結果を再確認できる性質です。
\end{tcolorbox}

\section{コマンド早見表}
\begin{longtable}{@{}L{59mm}p{89mm}@{}}
\toprule
コマンド & 目的\\
\midrule
\endfirsthead
\toprule コマンド & 目的\\\midrule\endhead
\code{bend version} & バージョン表示。\\
\code{bend guide} & インストール版に一致したガイド。\\
\code{bend base} & 標準ライブラリの参照。\\
\code{bend base Array} & 配列周辺のAPIを個別に調査。\\

\code{bend file.bend --check-only} & 型・停止性などのチェックだけ実施。\\
\code{bend file.bend} & チェックと実行／正規化。\\
\code{bend file.bend -o program} & ネイティブバイナリの生成。\\



\code{python3 -m unittest discover -s tests} & Python参照系の自己テスト。\\
\bottomrule
\end{longtable}
\section{主要資料}
\begin{itemize}
\item \url{https://github.com/bendlang/bend}\newline 公式リポジトリ、README、guide、Base、リリース。
\item \url{https://github.com/bendlang/bend/releases}（公式の最新版・リリース日時）。
\item \url{https://bend2.dev/}（非公式だが例題の豊富な学習リファレンス。公式\code{bend guide}が優先）。
\item \url{https://nohzafk.github.io/bend2-from-zero/}（Bend2の失敗例とライフゲームを含む独立教材。実装上の参考情報）。
\end{itemize}
\section{検証ログの書式}
\begin{lstlisting}[language={},numbers=none,caption={VERSION\_LOG.mdの記入例（数値は実測値に置換）}]
Date: YYYY-MM-DD
OS / Arch:
Bend version:
Bend native clang:
Commands:
  bend src/08_life_terminal.bend --check-only
  python3 scripts/check_bend.py
Result: PASS / FAIL / SKIP
Error output (if any):
Changes made after checking:
\end{lstlisting}
\section{最後の総合チェック}
\begin{enumerate}
\item \textbf{データ型}: \code{U32}、\code{Nat}、\code{String}、ユーザー定義型の違いを説明できます。
\item \textbf{関数}: 引数・返り値と\code{match}による場合分けを書けます。
\item \textbf{データ構造}: 配列の要素を読み書きし、所有権を戻り値で受け渡せます。
\item \textbf{関数型の考え方}: 高階関数、クロージャ、\code{Maybe}による失敗の扱いを説明できます。
\item \textbf{ライフゲーム}: 生死規則と同期更新を実装し、動作をテストで確かめられます。
\end{enumerate}

\backmatter
\end{document}
